\section{The naive cotangent complex}
\label{section-netherlander}

\noindent
Let $R \to S$ be a ring map. Denote $R[S]$ the polynomial ring
whose variables are the elements $s \in S$. Let's denote $[s] \in R[S]$
the variable corresponding to $s \in S$. Thus $R[S]$ is a free
$R$-module on the basis elements $[s_1] \ldots [s_n]$ where
$s_1, \ldots, s_n$ ranges over all unordered sequences of elements of $S$.
There is a canonical surjection
\begin{equation}
\label{equation-canonical-presentation}
R[S] \longrightarrow S,\quad [s] \longmapsto s
\end{equation}
whose kernel we denote $I \subset R[S]$. It is a simple observation that
$I$ is generated by the elements
$[s + s'] - [s] - [s']$, $[s][s'] - [ss']$ and $[r] - r$.
According to Lemma \ref{lemma-differential-seq}
there is a canonical map
\begin{equation}
\label{equation-naive-cotangent-complex}
I/I^2 \longrightarrow \Omega_{R[S]/R} \otimes_{R[S]} S
\end{equation}
whose cokernel is canonically isomorphic to $\Omega_{S/R}$. Observe that
the $S$-module $\Omega_{R[S]/R} \otimes_{R[S]} S$ is free on the generators
$\text{d}[s]$.

\begin{definition}
\label{definition-naive-cotangent-complex}
Let $R \to S$ be a ring map. The {\it naive cotangent complex}
$\NL_{S/R}$ is the chain complex (\ref{equation-naive-cotangent-complex})
$$
\NL_{S/R} = \left(I/I^2 \longrightarrow \Omega_{R[S]/R} \otimes_{R[S]} S\right)
$$
with $I/I^2$ placed in (homological) degree $1$ and
$\Omega_{R[S]/R} \otimes_{R[S]} S$ placed in degree $0$. We will denote
$H_1(L_{S/R}) = H_1(\NL_{S/R})$\footnote{This module is sometimes
denoted $\Gamma_{S/R}$ in the literature.} the homology in degree $1$.
\end{definition}

\noindent
Before we continue let us say a few words about the actual cotangent
complex (Cotangent, Section \ref{cotangent-section-cotangent-ring-map}).
Given a ring map $R \to S$ there exists a canonical simplicial
$R$-algebra $P_\bullet$ whose terms are polynomial algebras and
which comes equipped with a canonical homotopy equivalence
$$
P_\bullet \longrightarrow S
$$
The cotangent complex $L_{S/R}$ of $S$ over $R$ is defined as the chain
complex associated to the simplicial module
$$
\Omega_{P_\bullet/R} \otimes_{P_\bullet} S
$$
The naive cotangent complex as defined above is canonically isomorphic to
the truncation $\tau_{\leq 1}L_{S/R}$ (see
Homology, Section \ref{homology-section-truncations} and
Cotangent, Section \ref{cotangent-section-surjections}). In particular, it is
indeed the case that $H_1(\NL_{S/R}) = H_1(L_{S/R})$ so our definition is
compatible with the one using the cotangent complex. Moreover,
$H_0(L_{S/R}) = H_0(\NL_{S/R}) = \Omega_{S/R}$ as we've seen above.

\medskip\noindent
Let $R \to S$ be a ring map. A {\it presentation of $S$ over $R$} is
a surjection $\alpha : P \to S$ of $R$-algebras where $P$ is a polynomial
algebra (on a set of variables). Often, when $S$ is of finite type over $R$
we will indicate this by saying: ``Let $R[x_1, \ldots, x_n] \to S$ be a
presentation of $S/R$'', or
``Let $0 \to I \to R[x_1, \ldots, x_n] \to S \to 0$ be a presentation
of $S/R$'' if we want to indicate that $I$ is the kernel of the presentation.
Note that the map $R[S] \to S$ used to define the naive cotangent complex
is an example of a presentation.

\medskip\noindent
Note that for every presentation $\alpha$ we obtain a two term
chain complex of $S$-modules
$$
\NL(\alpha) : I/I^2 \longrightarrow \Omega_{P/R} \otimes_P S.
$$
Here the term $I/I^2$ is placed in degree $1$ and the term
$\Omega_{P/R} \otimes S$ is placed in degree $0$. The class of $f \in I$
in $I/I^2$ is mapped to $\text{d}f \otimes 1$ in $\Omega_{P/R} \otimes S$.
The cokernel of this complex is canonically $\Omega_{S/R}$,
see Lemma \ref{lemma-differential-seq}. We call the complex $\NL(\alpha)$
the {\it naive cotangent complex associated to the
presentation $\alpha : P \to S$ of $S/R$}. Note that if $P = R[S]$
with its canonical surjection onto $S$, then we recover $\NL_{S/R}$.
If $P = R[x_1, \ldots, x_n]$ then will sometimes use the notation
$I/I^2 \to \bigoplus_{i = 1, \ldots, n} S\text{d}x_i$
to denote this complex.

\medskip\noindent
Suppose we are given a commutative diagram
\begin{equation}
\label{equation-functoriality-NL}
\vcenter{
\xymatrix{
S \ar[r]_{\phi} & S' \\
R \ar[r] \ar[u] & R' \ar[u]
}
}
\end{equation}
of rings. Let $\alpha : P \to S$ be a presentation of $S$ over $R$ and let
$\alpha' : P' \to S'$ be a presentation of $S'$ over $R'$.
A {\it morphism of presentations from $\alpha : P \to S$ to
$\alpha' : P' \to S'$} is defined to be an $R$-algebra
map
$$
\varphi : P \to P'
$$
such that $\phi \circ \alpha = \alpha' \circ \varphi$. Note that
in this case $\varphi(I) \subset I'$, where $I = \Ker(\alpha)$
and $I' = \Ker(\alpha')$. Thus $\varphi$ induces a map
of $S$-modules $I/I^2 \to I'/(I')^2$ and by functoriality of
differentials also an $S$-module map
$\Omega_{P/R} \otimes S \to \Omega_{P'/R'} \otimes S'$.
These maps are compatible with the differentials of $\NL(\alpha)$ and
$\NL(\alpha')$ and we obtain a map of naive cotangent complexes
$$
\NL(\alpha) \longrightarrow \NL(\alpha').
$$
It is often convenient to consider the induced map
$\NL(\alpha) \otimes_S S' \to \NL(\alpha')$.

\medskip\noindent
In the special case that $P = R[S]$ and $P' = R'[S']$ the map
$\phi : S \to S'$ induces a canonical ring map
$\varphi : P \to P'$ by the rule $[s] \mapsto [\phi(s)]$.
Hence the construction above determines canonical(!) maps
of chain complexes
$$
\NL_{S/R} \longrightarrow \NL_{S'/R'},\quad\text{and}\quad
\NL_{S/R} \otimes_S S' \longrightarrow \NL_{S'/R'}
$$
associated to the diagram (\ref{equation-functoriality-NL}). Note that
this construction is compatible with composition: given a commutative
diagram
$$
\xymatrix{
S \ar[r]_{\phi} & S' \ar[r]_{\phi'} & S'' \\
R \ar[r] \ar[u] & R' \ar[u] \ar[r] & R'' \ar[u]
}
$$
we see that the composition of
$$
\NL_{S/R} \longrightarrow \NL_{S'/R'} \longrightarrow \NL_{S''/R''}
$$
is the map $\NL_{S/R} \to \NL_{S''/R''}$ given by the outer square.

\medskip\noindent
It turns out that $\NL(\alpha)$ is homotopy equivalent to $\NL_{S/R}$
and that the maps constructed above are well defined up to homotopy
(homotopies of maps of complexes are discussed in
Homology, Section \ref{homology-section-complexes}
but we also spell out the exact meaning of the statements in the lemma
below in its proof).

\begin{lemma}
\label{lemma-NL-homotopy}
Suppose given a diagram (\ref{equation-functoriality-NL}).
Let $\alpha : P \to S$ and $\alpha' : P' \to S'$ be presentations.
\begin{enumerate}
\item There exists a morphism of presentations from $\alpha$ to $\alpha'$.
\item Any two morphisms of presentations induce homotopic
morphisms of complexes $\NL(\alpha) \to \NL(\alpha')$.
\item The construction is compatible with compositions of morphisms
of presentations (see proof for exact statement).
\item If $R \to R'$ and $S \to S'$ are isomorphisms, then
for any map $\varphi$ of presentations from $\alpha$ to $\alpha'$
the induced map $\NL(\alpha) \to \NL(\alpha')$ is a homotopy equivalence
and a quasi-isomorphism.
\end{enumerate}
In particular, comparing $\alpha$ to the canonical presentation
(\ref{equation-canonical-presentation}) we conclude there is a
quasi-isomorphism $\NL(\alpha) \to \NL_{S/R}$ well defined
up to homotopy and compatible with all functorialities (up to homotopy).
\end{lemma}

\begin{proof}
Since $P$ is a polynomial algebra over $R$ we can write
$P = R[x_a, a \in A]$ for some set $A$.
As $\alpha'$ is surjective, we can choose
for every $a \in A$ an element $f_a \in P'$
such that $\alpha'(f_a) = \phi(\alpha(x_a))$. Let
$\varphi : P = R[x_a, a \in A] \to P'$ be the
unique $R$-algebra map such that $\varphi(x_a) = f_a$.
This gives the morphism in (1).

\medskip\noindent
Let $\varphi$ and $\varphi'$ morphisms of presentations from $\alpha$
to $\alpha'$. Let $I = \Ker(\alpha)$ and $I' = \Ker(\alpha')$.
We have to construct the diagonal map $h$ in the diagram
$$
\xymatrix{
I/I^2 \ar[r]^-{\text{d}}
\ar@<1ex>[d]^{\varphi'_1} \ar@<-1ex>[d]_{\varphi_1}
&
\Omega_{P/R} \otimes_P S
\ar@<1ex>[d]^{\varphi'_0} \ar@<-1ex>[d]_{\varphi_0}
\ar[ld]_h
\\
I'/(I')^2 \ar[r]^-{\text{d}}
&
\Omega_{P'/R'} \otimes_{P'} S'
}
$$
where the vertical maps are induced by $\varphi$, $\varphi'$ such that
$$
\varphi_1 - \varphi'_1 = h \circ \text{d}
\quad\text{and}\quad
\varphi_0 - \varphi'_0 = \text{d} \circ h
$$
Consider the map $\varphi - \varphi' : P \to P'$. Since both $\varphi$
and $\varphi'$ are compatible with $\alpha$ and $\alpha'$ we obtain
$\varphi - \varphi' : P \to I'$. This implies that
$\varphi, \varphi' : P \to P'$ induce the same $P$-module structure
on $I'/(I')^2$, since
$\varphi(p)i' - \varphi'(p)i' = (\varphi - \varphi')(p)i' \in (I')^2$.
Also $\varphi - \varphi'$ is $R$-linear and
$$
(\varphi - \varphi')(fg) =
\varphi(f)(\varphi - \varphi')(g) + (\varphi - \varphi')(f)\varphi'(g)
$$
Hence the induced map $D : P \to I'/(I')^2$ is a $R$-derivation.
Thus we obtain a canonical map $h : \Omega_{P/R} \otimes_P S \to I'/(I')^2$
such that $D = h \circ \text{d}$. 
A calculation (omitted) shows that $h$ is the desired homotopy.

\medskip\noindent
Suppose that we have a commutative diagram
$$
\xymatrix{
S \ar[r]_{\phi} & S' \ar[r]_{\phi'} & S'' \\
R \ar[r] \ar[u] & R' \ar[u] \ar[r] & R'' \ar[u]
}
$$
and that
\begin{enumerate}
\item $\alpha : P \to S$,
\item $\alpha' : P' \to S'$, and
\item $\alpha'' : P'' \to S''$
\end{enumerate}
are presentations. Suppose that
\begin{enumerate}
\item $\varphi : P \to P'$ is a morphism of presentations from
$\alpha$ to $\alpha'$ and
\item $\varphi' : P' \to P''$
is a morphism of presentations from $\alpha'$ to $\alpha''$.
\end{enumerate}
Then it is immediate that
$\varphi' \circ \varphi : P \to P''$
is a morphism of presentations from $\alpha$ to $\alpha''$ and that
the induced map $\NL(\alpha) \to \NL(\alpha'')$ of naive cotangent complexes
is the composition of the maps $\NL(\alpha) \to \NL(\alpha')$ and
$\NL(\alpha') \to \NL(\alpha'')$ induced by $\varphi$ and $\varphi'$.

\medskip\noindent
In the simple case of complexes with 2 terms a quasi-isomorphism
is just a map that induces an isomorphism on both the cokernel
and the kernel of the maps between the terms. Note that homotopic
maps of 2 term complexes (as explained above) define the same maps on
kernel and cokernel. Hence if $\varphi$ is a map from a presentation
$\alpha$ of $S$ over $R$ to itself, then the induced map
$\NL(\alpha) \to \NL(\alpha)$ is a quasi-isomorphism being homotopic
to the identity by part (2). To prove (4) in full generality, consider
a morphism $\varphi'$ from $\alpha'$ to $\alpha$ which exists by (1).
The compositions $\NL(\alpha) \to \NL(\alpha') \to \NL(\alpha)$ and
$\NL(\alpha') \to \NL(\alpha) \to \NL(\alpha')$ are homotopic to the identity
maps by (3), hence these maps are homotopy equivalences by definition.
It follows formally that both maps
$\NL(\alpha) \to \NL(\alpha')$ and $\NL(\alpha') \to \NL(\alpha)$ are
quasi-isomorphisms. Some details omitted.
\end{proof}

\begin{lemma}
\label{lemma-NL-polynomial-algebra}
Let $A \to B$ be a polynomial algebra. Then $\NL_{B/A}$ is homotopy equivalent
to the chain complex $(0 \to \Omega_{B/A})$ with $\Omega_{B/A}$
in degree $0$.
\end{lemma}

\begin{proof}
Follows from Lemma \ref{lemma-NL-homotopy}
and the fact that $\text{id}_B : B \to B$ is a presentation of $B$ over $A$
with zero kernel.
\end{proof}

\noindent
The following lemma is part of the motivation for introducing the
naive cotangent complex. The cotangent complex extends this
to a genuine long exact cohomology sequence. If $B \to C$ is a
local complete intersection, then one can extend the sequence with a
zero on the left, see
More on Algebra, Lemma \ref{more-algebra-lemma-transitive-lci-at-end}.

\begin{lemma}[Jacobi-Zariski sequence]
\label{lemma-exact-sequence-NL}
Let $A \to B \to C$ be ring maps. Choose a presentation
$\alpha : A[x_s, s \in S] \to B$ with kernel $I$. Choose a presentation
$\beta : B[y_t, t \in T] \to C$ with kernel $J$. Let
$\gamma : A[x_s, y_t] \to C$ be the induced presentation of $C$ with kernel
$K$. Then we get a canonical commutative diagram
$$
\xymatrix{
0 \ar[r] &
\Omega_{A[x_s]/A} \otimes C \ar[r] &
\Omega_{A[x_s, y_t]/A} \otimes C \ar[r] &
\Omega_{B[y_t]/B} \otimes C \ar[r] &
0 \\
&
I/I^2 \otimes C \ar[r] \ar[u] &
K/K^2 \ar[r] \ar[u] &
J/J^2 \ar[r] \ar[u] &
0
}
$$
with exact rows. We get the following exact sequence
of homology groups
$$
H_1(\NL_{B/A} \otimes_B C) \to
H_1(L_{C/A}) \to
H_1(L_{C/B}) \to
C \otimes_B \Omega_{B/A} \to
\Omega_{C/A} \to
\Omega_{C/B} \to 0
$$
of $C$-modules extending the sequence of
Lemma \ref{lemma-exact-sequence-differentials}.
If $\text{Tor}_1^B(\Omega_{B/A}, C) = 0$ and
$\text{Tor}_2^B(\Omega_{B/A}, C) = 0$, then
$H_1(\NL_{B/A} \otimes_B C) = H_1(L_{B/A}) \otimes_B C$.
\end{lemma}

\begin{proof}
The precise definition of the maps is omitted.
The exactness of the top row follows as the $\text{d}x_s$,
$\text{d}y_t$ form a basis for the middle module.
The map $\gamma$ factors
$$
A[x_s, y_t] \to B[y_t] \to C
$$
with surjective first arrow and second arrow equal to $\beta$.
Thus we see that $K \to J$ is surjective.
Moreover, the kernel of the first displayed arrow is
$IA[x_s, y_t]$. Hence $I/I^2 \otimes C$ surjects onto the
kernel of $K/K^2 \to J/J^2$. Finally, we can use
Lemma \ref{lemma-NL-homotopy}
to identify the terms as homology groups of the naive
cotangent complexes.

\medskip\noindent
The final assertion is a statement in homological algebra.
Recall that $\NL_{B/A} = (N^{-1} \to N^0)$ is a two term
complex of $B$-modules with $N^0$ free and cohomology modules
$H^0 = \Omega_{B/A}$ and $H^{-1} = H_1(L_{B/A})$.
Write $M \subset N^0$ for the image of the differential.
If $\text{Tor}_1^B(H^0, C) = 0$, then we have an exact sequence
$$
0 \to M \otimes_B C \to N^0 \otimes_B C \to H^0 \otimes_B C \to 0
$$
Since $N^0$ is free, we also see that $\text{Tor}_2^B(H^0, C) =
\text{Tor}_1^B(M, C)$. Hence if $\text{Tor}_2^B(H^0, C) = 0$
then we also have an exact sequence
$$
0 \to H^{-1} \otimes_B C \to N^{-1} \otimes_B C \to M \otimes_B C \to 0
$$
Putting everything together we see that if
$\text{Tor}_1^B(H^0, C) = 0$ and $\text{Tor}_2^B(H^0, C) = 0$,
then $H^{-1} \otimes_B C$ is the kernel of
$N^{-1} \otimes_B C \to N^0 \otimes_B C$ as desired.
\end{proof}

\begin{remark}
\label{remark-composition-homotopy-equivalent-to-zero}
Let $A \to B$ and $\phi : B \to C$ be ring maps.
Then the composition $\NL_{B/A} \to \NL_{C/A} \to \NL_{C/B}$ is
homotopy equivalent to zero. Namely, this composition is the functoriality
of the naive cotangent complex for the square
$$
\xymatrix{
B \ar[r]_\phi & C \\
A \ar[r] \ar[u] & B \ar[u]
}
$$
Write $J = \Ker(B[C] \to C)$. An explicit homotopy is given by the map
$\Omega_{A[B]/A} \otimes_{A[B]} B \to J/J^2$ which maps the basis element
$\text{d}[b]$ to the class of $[\phi(b)] - b$ in $J/J^2$.
\end{remark}

\begin{lemma}
\label{lemma-NL-surjection}
Let $A \to B$ be a surjective ring map with kernel $I$.
Then $\NL_{B/A}$ is homotopy equivalent to the chain complex
$(I/I^2 \to 0)$ with $I/I^2$ in degree $1$. In particular
$H_1(L_{B/A}) = I/I^2$.
\end{lemma}

\begin{proof}
Follows from Lemma \ref{lemma-NL-homotopy}
and the fact that $A \to B$ is a presentation of $B$ over $A$.
\end{proof}

\begin{lemma}
\label{lemma-application-NL}
Let $A \to B \to C$ be ring maps. Assume $A \to C$ is surjective (so
also $B \to C$ is). Denote $I = \Ker(A \to C)$ and
$J = \Ker(B \to C)$. Then the sequence
$$
I/I^2 \to J/J^2 \to \Omega_{B/A} \otimes_B B/J \to 0
$$
is exact.
\end{lemma}

\begin{proof}
Follows from
Lemma \ref{lemma-exact-sequence-NL}
and the description of the naive cotangent complexes
$\NL_{C/B}$ and $\NL_{C/A}$ in Lemma \ref{lemma-NL-surjection}.
\end{proof}

\begin{lemma}[Flat base change]
\label{lemma-change-base-NL}
Let $R \to S$ be a ring map. Let $\alpha : P \to S$ be a presentation.
Let $R \to R'$ be a flat ring map.
Let $\alpha' : P \otimes_R R' \to S' = S \otimes_R R'$
be the induced presentation.
Then $\NL(\alpha) \otimes_R R' = \NL(\alpha) \otimes_S S' = \NL(\alpha')$.
In particular, the canonical map
$$
\NL_{S/R} \otimes_S S' \longrightarrow \NL_{S \otimes_R R'/R'}
$$
is a homotopy equivalence if $R \to R'$ is flat.
\end{lemma}

\begin{proof}
This is true because
$\Ker(\alpha') = R' \otimes_R \Ker(\alpha)$
since $R \to R'$ is flat.
\end{proof}

\begin{lemma}
\label{lemma-colimits-NL}
Let $R_\lambda \to S_\lambda$ be a system of ring maps over
the directed set $\Lambda$.
Set $R = \colim R_\lambda$ and $S = \colim S_\lambda$.
Then $\NL_{S/R} = \colim \NL_{S_\lambda/R_\lambda}$.
\end{lemma}

\begin{proof}
Recall that $\NL_{S/R}$ is the complex
$I/I^2 \to \bigoplus_{s \in S} S\text{d}[s]$ where $I \subset R[S]$
is the kernel of the canonical presentation $R[S] \to S$.
Now it is clear that $R[S] = \colim R_\lambda[S_\lambda]$ and similarly
that $I = \colim I_\lambda$ where
$I_\lambda = \Ker(R_\lambda[S_\lambda] \to S_\lambda)$.
Hence the lemma is clear.
\end{proof}

\begin{lemma}
\label{lemma-NL-of-localization}
If $S \subset A$ is a multiplicative subset of $A$, then
$\NL_{S^{-1}A/A}$ is homotopy equivalent to the zero complex.
\end{lemma}

\begin{proof}
Since $A \to S^{-1}A$ is flat we see that
$\NL_{S^{-1}A/A} \otimes_A S^{-1}A \to \NL_{S^{-1}A/S^{-1}A}$
is a homotopy equivalence by flat base change
(Lemma \ref{lemma-change-base-NL}). Since the source of the arrow
is isomorphic to $\NL_{S^{-1}A/A}$ and the target of the arrow is
zero (by Lemma \ref{lemma-NL-surjection}) we win.
\end{proof}

\begin{lemma}
\label{lemma-NL-localize-bottom}
Let $S \subset A$ is a multiplicative subset of $A$.
Let $S^{-1}A \to B$ be a ring map.
Then $\NL_{B/A} \to \NL_{B/S^{-1}A}$ is a homotopy equivalence.
\end{lemma}

\begin{proof}
Choose a presentation $\alpha : P \to B$ of $B$ over $A$.
Then $\beta : S^{-1}P \to B$ is a presentation of $B$ over $S^{-1}A$.
A direct computation shows that we have $\NL(\alpha) = \NL(\beta)$
which proves the lemma as the naive cotangent complex is well defined
up to homotopy by Lemma \ref{lemma-NL-homotopy}.
\end{proof}

\begin{lemma}
\label{lemma-principal-localization-NL}
\begin{slogan}
The formation of the naive cotangent complex commutes with localization
at an element.
\end{slogan}
Let $A \to B$ be a ring map. Let $g \in B$. Suppose $\alpha : P \to B$
is a presentation with kernel $I$. Then a presentation of $B_g$ over $A$ is
the map
$$
\beta : P[x] \longrightarrow B_g
$$
extending $\alpha$ and sending $x$ to $1/g$.
The kernel $J$ of $\beta$ is generated by $I$ and the element $f x - 1$
where $f \in P$ is an element mapped to $g \in B$ by $\alpha$. In this
situation we have
\begin{enumerate}
\item $J/J^2 = (I/I^2)_g \oplus B_g (f x - 1)$,
\item $\Omega_{P[x]/A} \otimes_{P[x]} B_g =
\Omega_{P/A} \otimes_P B_g \oplus B_g \text{d}x$,
\item $\NL(\beta) \cong
\NL(\alpha) \otimes_B B_g \oplus (B_g \xrightarrow{g} B_g)$
\end{enumerate}
Hence the canonical map $\NL_{B/A} \otimes_B B_g \to \NL_{B_g/A}$
is a homotopy equivalence.
\end{lemma}

\begin{proof}
Since $P[x]/(I, fx - 1) = B[x]/(gx - 1) = B_g$ we get the statement about
$I$ and $fx - 1$ generating $J$. Consider the commutative diagram
$$
\xymatrix{
0 \ar[r] &
\Omega_{P/A} \otimes B_g \ar[r] &
\Omega_{P[x]/A} \otimes B_g \ar[r] &
\Omega_{B[x]/B} \otimes B_g \ar[r] &
0 \\
&
(I/I^2)_g \ar[r] \ar[u] &
J/J^2 \ar[r] \ar[u] &
(gx - 1)/(gx - 1)^2 \ar[r] \ar[u] &
0
}
$$
with exact rows of Lemma \ref{lemma-exact-sequence-NL}.
The $B_g$-module $\Omega_{B[x]/B} \otimes B_g$ is free of
rank $1$ on $\text{d}x$. The element $\text{d}x$ in the
$B_g$-module $\Omega_{P[x]/A} \otimes B_g$ provides
a splitting for the top row. The element $gx - 1 \in (gx - 1)/(gx - 1)^2$
is mapped to $g\text{d}x$ in $\Omega_{B[x]/B} \otimes B_g$
and hence $(gx - 1)/(gx - 1)^2$ is free of rank $1$ over $B_g$.
(This can also be seen by arguing that $gx - 1$ is a nonzerodivisor
in $B[x]$ because it is a polynomial with invertible constant term
and any nonzerodivisor gives a quasi-regular sequence of length $1$
by Lemma \ref{lemma-regular-quasi-regular}.)

\medskip\noindent
Let us prove $(I/I^2)_g \to J/J^2$ injective. Consider the $P$-algebra map
$$
\pi : P[x] \to (P/I^2)_f = P_f/I_f^2
$$
sending $x$ to $1/f$. Since $J$ is generated by $I$ and $fx - 1$
we see that $\pi(J) \subset (I/I^2)_f = (I/I^2)_g$. Since this
is an ideal of square zero we see that $\pi(J^2) = 0$.
If $a \in I$ maps to an element of $J^2$ in $J$, then
$\pi(a) = 0$, which implies that $a$ maps to zero in $I_f/I_f^2$.
This proves the desired injectivity.

\medskip\noindent
Thus we have a short exact sequence of two term complexes
$$
0 \to \NL(\alpha) \otimes_B B_g \to \NL(\beta)
\to (B_g \xrightarrow{g} B_g) \to 0
$$
Such a short exact sequence can always be split in the category of
complexes. In our particular case we can take as splittings
$$
J/J^2 = (I/I^2)_g \oplus B_g (fx - 1)\quad\text{and}\quad
\Omega_{P[x]/A} \otimes B_g = \Omega_{P/A} \otimes B_g \oplus
B_g (g^{-2}\text{d}f + \text{d}x)
$$
This works because
$\text{d}(fx - 1) = x\text{d}f + f \text{d}x =
g(g^{-2}\text{d}f + \text{d}x)$
in $\Omega_{P[x]/A} \otimes B_g$.
\end{proof}

\begin{lemma}
\label{lemma-localize-NL}
Let $A \to B$ be a ring map. Let $S \subset B$ be a multiplicative subset.
The canonical map $\NL_{B/A} \otimes_B S^{-1}B \to \NL_{S^{-1}B/A}$
is a quasi-isomorphism.
\end{lemma}

\begin{proof}
We have $S^{-1}B = \colim_{g \in S} B_g$ where we think of $S$
as a directed set (ordering by divisibility), see
Lemma \ref{lemma-localization-colimit}.
By Lemma \ref{lemma-principal-localization-NL} each of the maps
$\NL_{B/A} \otimes_B B_g \to \NL_{B_g/A}$
are quasi-isomorphisms.
The lemma follows from Lemma \ref{lemma-colimits-NL}.
\end{proof}

\begin{lemma}
\label{lemma-sum-two-terms}
Let $R$ be a ring.
Let $A_1 \to A_0$, and $B_1 \to B_0$ be
two term complexes. Suppose that there exist
morphisms of complexes $\varphi : A_\bullet \to B_\bullet$
and $\psi : B_\bullet \to A_\bullet$ such that
$\varphi \circ \psi$ and $\psi \circ \varphi$ are
homotopic to the identity maps.
Then $A_1 \oplus B_0 \cong B_1 \oplus A_0$ as
$R$-modules.
\end{lemma}

\begin{proof}
Choose a map $h : A_0 \to A_1$ such that
$$
\text{id}_{A_1} - \psi_1 \circ \varphi_1 = h \circ d_A
\text{ and }
\text{id}_{A_0} - \psi_0 \circ \varphi_0 = d_A \circ h.
$$
Similarly, choose a map $h' : B_0 \to B_1$ such that
$$
\text{id}_{B_1} - \varphi_1 \circ \psi_1 = h' \circ d_B
\text{ and }
\text{id}_{B_0} - \varphi_0 \circ \psi_0 = d_B \circ h'.
$$
A trivial computation shows that
$$
\left(
\begin{matrix}
\text{id}_{A_1} & -h' \circ \psi_1 + h \circ \psi_0 \\
0 & \text{id}_{B_0}
\end{matrix}
\right)
=
\left(
\begin{matrix}
\psi_1 & h \\
-d_B & \varphi_0
\end{matrix}
\right)
\left(
\begin{matrix}
\varphi_1 & - h' \\
d_A & \psi_0
\end{matrix}
\right)
$$
This shows that both matrices on the right hand side
are invertible and proves the lemma.
\end{proof}

\begin{lemma}
\label{lemma-conormal-module}
Let $R \to S$ be a ring map of finite type.
For any presentations $\alpha : R[x_1, \ldots, x_n] \to S$, and
$\beta : R[y_1, \ldots, y_m] \to S$ we have
$$
I/I^2 \oplus S^{\oplus m} \cong J/J^2 \oplus S^{\oplus n}
$$
as $S$-modules where $I = \Ker(\alpha)$ and $J = \Ker(\beta)$.
\end{lemma}

\begin{proof}
See Lemmas \ref{lemma-NL-homotopy} and \ref{lemma-sum-two-terms}.
\end{proof}

\begin{lemma}
\label{lemma-conormal-module-localize}
Let $R \to S$ be a ring map of finite type.
Let $g \in S$. For any presentations
$\alpha : R[x_1, \ldots, x_n] \to S$, and
$\beta : R[y_1, \ldots, y_m] \to S_g$ we have
$$
(I/I^2)_g \oplus S^{\oplus m}_g \cong J/J^2 \oplus S_g^{\oplus n}
$$
as $S_g$-modules where
$I = \Ker(\alpha)$ and $J = \Ker(\beta)$.
\end{lemma}

\begin{proof}
Let $\beta' : R[x_1, \ldots, x_n, x] \to S_g$ be the
presentation of Lemma \ref{lemma-principal-localization-NL}
constructed starting with $\alpha$. Then we know that
$\NL(\alpha) \otimes_S S_g$ is homotopy equivalent to $\NL(\beta')$.
We know that $\NL(\beta)$ and $\NL(\beta')$ are homotopy equivalent
by Lemma \ref{lemma-NL-homotopy}. We conclude that
$\NL(\alpha) \otimes_S S_g$ is homotopy equivalent to $\NL(\beta)$.
Finally, we apply Lemma \ref{lemma-conormal-module}.
\end{proof}







